% !TEX root = ../b-rg.tex
\chapter{Coordination and Subordination}\label{ch:complex-sentences}

This chapter describes how clauses and phrases are combined: coordination and the ellipsis it allows, complement clauses, adverbial clauses, conditional sentences, participial clauses and reported speech. The choice between indicative and subjunctive in subordinate clauses is treated in Section~\ref{sec:mood}, relative clauses in Chapter~\ref{ch:relative-clauses}, and indirect questions in Section~\ref{sec:indirect-questions}.

All subordinate clauses have the subject before the verb, and a fronted subordinate clause does not cause inversion in the main clause that follows it (Section~\ref{sec:v2}).

\section{Coordination}\label{sec:coordination}
The coordinating conjunctions are \bl{e} `and', \bl{ma} `but' and \bl{ben} `or'. They stand between the conjuncts, and only one is used however many conjuncts there are (\bl{l'amand, agl, ugnlac eð arjon} `the almond, garlic, onion and pepper'). Before a vowel, \bl{e} takes the form \bl{eð}, as \bl{ou} `whether' takes \bl{ouð}. This includes words beginning with \bl{y} (\bl{eð y caðer} `and the chairs'), but not those beginning with \bl{h} (\bl{e Halina}).
    \glsen{Le no's ja ac yent, ma sy carc veyaçour es ac deja.}{le nɔz ʒa ak jɛnt ma si kark ˌvi.jaˈtsur ɛz ak deˈʒa}{{\Third\Sg.\M} {\Neg=be} yet here {arrive.\Pst.\Ptcp} but {\Third\Sg.\Poss} load travel be here already}{He hasn't arrived yet, but his luggage is already here.}
`Either \ldots{} or' is \bl{ou \ldots{} ben}, with \bl{ou} before the first conjunct and \bl{ben} before the second. `Neither \ldots{} nor' is \bl{necun \ldots{} ben} (Section~\ref{sec:neg-coordination}).
    \glsen{Ou tu part ben j'oc fay.}{u ti part bɛn ʒɔk fe}{either {\Second\Sg} leave or {\First\Sg=\Prox.\Sg} do}{Either you leave or I do.}
`As well as' is \bl{tan com}: \bl{Sy dou fraðr antir moroy com sugtour, tan com sy fraðr postir Peðr} `Her two older brothers died in infancy, as well as her younger brother Peter'.

Coordinated subjects take a plural verb (Section~\ref{sec:agreement}). When coordinated verbs share a subject, it is expressed only once, before the first verb, and each verb has its own object.
    \glsen{La tenoy sy veðrdon n'oc ort e soureyeu.}{la teˈnɔj si ˌvɛð.r̩ˈdɔn nɔk ɔrt e ˌsu.riˈjau}{{\Third\Sg.\F} {hold-\Pst} {\Third\Sg.\Poss} prize {in=\Prox.\Sg} rise and {smile-\Pst}}{She held up her trophy and grinned.}
The same holds for long series of verbs (\bl{Zoe descenoy y novel leyour peripaton quaucant e gasau y gið embrostað} `Zoe strode out of the new unmanned lifter and peered at the ornate signposts'), and for coordinated infinitives in instructions (\bl{Ec covauclir e cogr un hour e demeg sur foy preuç} `Place on a lid and cook for an hour and a half on a low fire').

Clauses are also linked by adverbs and by conjunction-like words that stand outside the clause. \bl{Soulor} `then, next' joins clauses without inverting the subject, unlike the adverb \bl{lor} `then' (Section~\ref{sec:v2}). \bl{Donc} `so' introduces a consequence (\bl{Inacavabr es soup a kigmau pascr con digntað, donc nel attendr} `It's impossible to eat lardon soup with dignity, so please don't try'). \bl{Malgrað} `however, though' and \bl{nentamen} `nonetheless' are adverbs and are usually set off by commas.

\section{Ellipsis}\label{sec:ellipsis}
A verb repeated in a second conjunct may be left out (gapping): \bl{Ana am pom, e Tom coðogn} `Ana likes apples, and Tom pears', and with \bl{plagr} `please', \bl{Pom plagn a Ana, e coðogn a Tom}.

A verb phrase is not left out on its own. It is replaced by the pro-verb \bl{oc fair} `do so', with \bl{oc} in the clitic position. The pro-verb takes the tense and mood of the verb it replaces (\bl{ben ou Marc oc faç} `or whether Marc does', with the subjunctive \bl{faç} after \bl{ou}). A modal or auxiliary keeps its own verb, with \bl{oc} as its object (\bl{Ja, j'oc sa} `Yes, I can'; Section~\ref{sec:answers}).
    \glsen{Mey, j'oc fey.}{mi ʒɔk fi}{{\First\Sg.\Dsj} {\First\Sg=\Prox.\Sg} {do.\Pst}}{I did.}
In a negative second conjunct, the subject may simply be followed by \bl{noc}: \bl{Marc e Halina venoirn, ma Theo noc} `Marc and Halina came, but Theo didn't', or with the pro-verb, \bl{ma Theo oc fey noc} (Section~\ref{sec:neg-coordination}). Alternatively, the contrasted subjects can be set out as disjunctive pronouns: \bl{Ou tey, tu part ben mey} `Either you leave or I do' (literally `either you, you leave, or me').

\section{Complement clauses}\label{sec:complement-clauses}

\subsection{Finite complements}
A finite complement clause is introduced by \bl{ig} `that', which has the euphonic by-form \bl{l'ig} after a word ending in a vowel. After verbs of saying, writing and thinking, \bl{ig} is often left out.
    \glsen{La scavoy la vole ac yembr je Mart.}{la xaˈvɔj la voˈle ak ˈjɛm.br̩ ʒe mart}{{\Third\Sg.\F} {write-\Pst} {\Third\Sg.\F} {will-\Ipfv} here {reach-\Inf} {day.of} Tuesday}{She wrote that she would arrive on Tuesday.}
The complement is in the indicative when its content is asserted, and in the subjunctive when it is wanted, planned, feared or doubted, or when a verb of knowing is negated or questioned (Section~\ref{sec:subjunctive}). A complement clause may also be the subject of a predicate. It then follows the predicate, usually with \bl{ig}, though not always (\bl{Rumorað es ig l'oç prinç sey zamicað an masmorrat} `That prince is rumoured to be locked in a dungeon', \bl{Incossirabr es nos neleu nossy vecin reconnoscr\ldots} `It is absurd that we shouldn't speak to our neighbour\ldots').

\subsection{Infinitival complements}
When the subject of the complement is the same as that of the main verb, the complement is usually an infinitive. This is the rule with verbs of wanting (\bl{Jo desir partir} `I want to leave', against \bl{Jo desir ig tu part} `I want you to leave'; Section~\ref{sec:subjunctive}), and it is common with verbs of saying and thinking (\bl{Ðe disrn empley broucar n'un foncçondry} `They said they had jobs at an agency'; Section~\ref{sec:tense-subordinate}). Some verbs take the infinitive with a preposition. \bl{Promettr} `promise' takes \bl{de}, while \bl{peðir} `ask' takes \bl{a}, with the person asked as an indirect object. \bl{Peðir} can also take a subjunctive clause (\bl{Jo peðeu (l'ig) l'attenois} `I asked that she wait'), but the infinitive is preferred.
    \glsen{Jo promis de scaur a casc semman.}{ʒo proˈmɪz de xor a kax sɛˈman}{{\First\Sg} {promise.\Pst} of {write-\Inf} at every week}{I promised to write every week.}
    \glsen{Jo luy peðeu a attendr.}{ʒo laj peˈðau a aˈtɛn.dr̩}{{\First\Sg} {\Third\Sg.\Obl} {ask-\Pst} to {wait-\Inf}}{I asked her to wait.}
As in every non-finite verb phrase, an object precedes the infinitive (Section~\ref{sec:basic-order}).

\subsection{Complements of perception verbs}
Verbs of perception take a bare infinitive, whether the action perceived is complete or ongoing. The person perceived is an object clitic on the perception verb, and the infinitive's own object precedes it.
    \glsen{Jo le vis y majon salir.}{ʒo le vɪz i maˈʒɔn saˈlɪr}{{\First\Sg} {\Third\Sg.\M.\Acc} {see.\Pst} {\Def} house {leave-\Inf}}{I saw him leave the house.}
    \glsen{Jo l'oyeu cantar ny cocin.}{ʒo lɔˈjau kanˈtar ni koˈtsɪn}{{\First\Sg} {\Third\Sg.\F.\Acc=hear-\Pst} {sing-\Inf} {in.\Def} kitchen}{I heard her singing in the kitchen.}
A progressive infinitive is also possible (\bl{Jo l'oyeu star cantant ny cocin} `I heard her be singing in the kitchen'), and so is \bl{oc} with a present participle, which adds a non-restrictive comment (\bl{Jo l'oyeu oc cantant\ldots} `I heard her, who was singing\ldots'; Section~\ref{sec:nonrestrictive}).

\section{Adverbial clauses}\label{sec:adverbial-clauses}
Adverbial clauses are introduced by a subordinator, which often contains \bl{ig}. The main subordinators, with the mood of their clause, are listed below.
\begin{table}[ht]
    \centering\small
    \begin{tabular}{l l l}
        \multicolumn{3}{c}{Adverbial subordinators} \\
        \hline
        meaning & subordinator & verb form \\
        \hline
        `when' & \bl{cant} & indicative; subjunctive of the future \\
        `when(ever)' & \bl{candon} & indicative \\
        `while' & \bl{tandic}, \bl{cant} & indicative \\
        `as soon as' & \bl{pront la} & indicative, or infinitive \\
        `before' & \bl{pre} & subjunctive, or infinitive \\
        `until' & \bl{fin a} & subjunctive; indicative of a past fact \\
        `after' & \bl{posc} & indicative, or infinitive \\
        `once' & \bl{alcant ig} & indicative \\
        `since' & \bl{des} & indicative \\
        `because' & \bl{dec} & indicative \\
        `so that' & \bl{a fin ig} & subjunctive \\
        `lest, so that \ldots{} not' & \bl{neðan} & subjunctive, or infinitive \\
        `in the hope that' & \bl{ne spien (ig)} & subjunctive \\
        `although' & \bl{degnant (ig)} & subjunctive \\
        `but, however' & \bl{con tan} & indicative \\
        `even though' & \bl{con tan} & subjunctive \\
        `unless' & \bl{preðauc} & subjunctive \\
        `if' & \bl{ða} & see Section~\ref{sec:conditionals} \\
        `as if' & \bl{com ða}, \bl{com ig} & indicative; subjunctive if unreal \\
        `without' & \bl{lon} & infinitive \\
        \hline
    \end{tabular}
\end{table}
Where an infinitive is possible, it is used when the subject is the same as that of the main clause. The same subordinators serve as prepositions before a noun phrase: \bl{fin a quelt} `until dinner', \bl{posc y var} `since the war', \bl{pre siecr ogtoç} `before the eighteenth century'.

\subsection{Time}
\bl{Cant} `when' takes the indicative of a past or habitual event, and the subjunctive of a future one (\bl{cant y rafayat revin} `when the tailor comes back'; Section~\ref{sec:sbjv-adverbial}). \bl{Candon} means `whenever, at times when' (\bl{Ðe amn queldar tras port candon cay bel} `They like to dine outside when it's fine out'), but is also used of a single past occasion (\bl{Candon j'entrau, y babe er dorment} `When I came in, the baby was sleeping').
    \glsen{J'attenoy cos infruniscent tandic y grimmour cave yert.}{ˌʒa.teˈnɔj kɔz ɪnˌfri.niˈxɛnt tanˈdɪk i grɪˈmur kaˈve jɛrt}{{\First\Sg=wait-\Pst} {\Adv} bored while {\Def} computer {get-\Ipfv} awake}{I waited impatiently while the computer booted up.}
\bl{Cant} also has a contrastive sense, `while, whereas': \bl{cant y dovem se dey parcanç ag fin -iscas} `while the second (element) may be from the suffix -iscas'.

\bl{Pront la} `as soon as' takes a finite clause, or an infinitive when the subject is the same.
    \glsen{Sceye un rasamar de conement pront la y varous cof yert.}{xiˈje ɪn ˌra.zaˈmar de ˌko.neˈmɛnt prɔnt la i vaˈruz kɔf jɛrt}{{there.be-\Ipfv} {\Indf} rush of customer immediate there {\Def} shop {get.\Pst} open}{There was a rush of customers as soon as the shop opened.}
    \glsen{Y tas smuïsceu pront la y paið quinar.}{i taz ˌsma.jiˈxaw prɔnt la i peθ kwiˈnar}{{\Def} saucer {shatter-\Pst} immediate there {\Def} floor {collide-\Inf}}{The saucer shattered as soon as it hit the floor.}
\bl{Pre} `before' takes the subjunctive, with no \bl{ig}, or an infinitive with the same subject (\bl{Pre decidr, vos pouð y tarn me loitar} `Before you decide, you can rent out the pub to me'). With a negative main clause it corresponds to English `until': \bl{Jo no poð lou affaçar pre vienn y flottisc} `I can't confront them until reinforcements arrive'.
    \glsen{Pre tu vay, closisc y fenestr.}{pre ti ve kloˈzɪx i feˈnɛs.tr̩}{before {\Second\Sg} {go.\Sbjv} close {\Def} window}{Before you go, close the window.}
\bl{Fin a} `until' (literally `end at') takes the subjunctive of something still to come (\bl{Torneyeð y cout fin a y zygom sey ovart} `Turn the crank until the valve's open'), but the indicative of a past fact.
    \glsen{J'asseye n'oc hour fin a tandem alandau my convey.}{ˌʒa.siˈje nɔk ur fɪn a tanˈdɛm ˌa.lanˈdo mi kɔnˈvi}{{\First\Sg=sit-\Ipfv} {in=\Prox.\Sg} time end at {at.last} {dock-\Pst} {\First\Sg.\Poss} train}{I sat for a bit until my train came at last.}
`After' is \bl{posc}, with a finite clause in any tense of the indicative, or with an infinitive when the subject is the same (\bl{Jo vil descombr posc l'oc capparol finir} `I'll go to bed after I finish this chapter'). `Since' is \bl{des} with a finite clause (\bl{Y molin produg alcun zucar des insorgent lo quariçtovaurn} `The mill hasn't produced any sugar since rebels sabotaged it'), or \bl{posc} or \bl{des} with a noun phrase (\bl{Le n'a bucet beut posc y var} `He hasn't smoked since the war').
An earlier event after \bl{posc} is in the pluperfect with the preterite auxiliary, so that the main event follows a completed one. \bl{Alcant ig} `once' instead takes the pluperfect with the imperfect auxiliary \bl{er}, so that the main event follows once the resulting state holds: \bl{Alcant ig le'r partið, \ldots} `Once he had left, \ldots'.
    \glsen{Posc le fo partið, nos ovroim y vin.}{pɔx le fo parˈtɪθ nɔz oˈvrɔjm i vɪn}{after {\Third\Sg.\M} {be.\Pst} {leave.\Pst.\Ptcp} {\First\Pl} {open-\Pst.\First\Pl} {\Def} wine}{After he had left, we opened the wine.}

\subsection{Cause}
`Because' is \bl{dec}, with the indicative. `Because of' with a noun phrase is \bl{desig} (\bl{Me val noc entrar desig mell'azig hortumesc} `I shouldn't come in because of my filthy shoes') or \bl{deut a}.
    \glsen{Stou ty sour convoccar dec la's tan ricottabr.}{stu ti sur ˌkɔn.vɔˈkar dɛk laz tan ˌri.kɔˈta.br̩}{{be.needed} {\Second\Sg.\Poss} sister {invite-\Inf} because {\Third\Sg.\F=be} so funny}{You should invite your sister; she's so funny.}
A \bl{dec} clause may also consist of a bare participle, with the subject understood from the main clause: \bl{Le vendra tart dec caut flaccið ne mot routoir} `He'll be late since he's stuck in a traffic jam'.

\subsection{Purpose}
`So that' is \bl{a fin ig} and `lest, so that \ldots{} not' is \bl{neðan}, both with the subjunctive. \bl{Neðan} takes an infinitive when the subject is the same (Section~\ref{sec:sbjv-adverbial}). A hoped-for result is expressed with \bl{ne spien (ig)} `in the hope (that)', likewise with the subjunctive, or with \bl{de} and an infinitive.
    \glsen{J'azia ne spien ig un hom interrompois my laur d'asn.}{ʒaˈzja ne spjɛn aj ɪn ɔm ɪnˌtɛ.rɔmˈpɔjz mi lor ˈda.sn̩}{{\First\Sg=stand-\Ipfv} in hope {\Comp} {\Indf} person {interrupt-\Ipfv.\Sbjv} {\First\Sg.\Poss} toil {of=donkey}}{I stood, hoping someone would interrupt the tedium.}
After a verb of motion, purpose is expressed by a bare infinitive (\bl{Me storra le vin y tabr despondr} `I need him to come and set the table') or by \bl{por} `for' with an infinitive (\bl{Y segr sera donc un vecin, ac yent por sbay dar ag forannour} `The man must be a local, come to gawk at the foreigners').

\subsection{Concession}
`Although' is \bl{degnant}, usually followed by \bl{ig}, with the subjunctive. In a past context this is the imperfect subjunctive, or the pluperfect subjunctive for an earlier event (\bl{Degnant le fos parsið con nis, le disceu ny scitaçon} `Although he had studied hard, he failed the exam'; Section~\ref{sec:pluperfect}).
    \glsen{Degnant ig cous tart, y varous remanen yert.}{dɪjˈnant aj kus tart i vaˈruz ˌre.maˈnɛn jɛrt}{although {\Comp} {fall.\Ipfv.\Sbjv} late {\Def} shop {remain-\Ipfv.\Third\Pl} open}{Although it was late, the shops were still open.}
\bl{Con tan} has two senses, distinguished by mood. With the indicative it is adversative, `but, nevertheless, however': \bl{L'ec gran accommission corrogationnal descrein y liesc romannesc, con tan l'ec dezein nolen l'oç stikettanç accarr pre un siecr may} `These great infrastructure projects characterise the Romantic period, although these decades would not gain that appellation for another century'. With the subjunctive it is concessive, `even though, despite that': \bl{Jo lo nommau runam com rouvaç, con tan, jo l'admet, lo fos y rouvaç ajoutandessem meyon jamay sbolognað} `I privately called it robbery, though I confess it was the pleasantest robbery I have ever suffered'. In this sense \bl{con tan} is the everyday word, and \bl{degnant ig} is more formal and written (\bl{La sorteu, degnant ig cous ploy}).
    \glsen{La sorteu con tan cous ploy.}{la sɔrˈtaw kɔn tan kus plɔj}{{\Third\Sg.\F} {go.out-\Pst} with {so.much} {fall.\Ipfv.\Sbjv} rain}{She went out, even though it was raining.}

\subsection{Manner and comparison}
`As if, as though' is \bl{com ða} or \bl{com ig}, with the indicative as a rule (\bl{Ðe ayen heu lastumnesc, com ig ðe no'rn for pog dormið} `They looked haggard, as though they hadn't slept much'). The present subjunctive marks the comparison as contrary to fact (\bl{Le parol com ða le saif tout} `He talks as if he knew everything (but he doesn't)'; Section~\ref{sec:sbjv-adverbial}). Comparative clauses with \bl{a} `than' and \bl{com} `as' are described in Chapter~\ref{ch:comparison}.

\subsection{Without}
`Without doing' is \bl{lon} with an infinitive, whose object precedes it.
    \glsen{La parteu lon un parol digr.}{la parˈtaw lɔn ɪn paˈrɔl ˈdaj.r̩}{{\Third\Sg.\F} {leave-\Pst} without {\Indf} word {say-\Inf}}{She left without saying a word.}

\section{Conditional sentences}\label{sec:conditionals}
Conditional clauses are introduced by \bl{ða} `if'. The tenses of the two clauses combine as described in Section~\ref{sec:conditional-sentences}:
\begin{itemize}
    \item an \textbf{open condition} has the present after \bl{ða}, and a future, present or imperative in the main clause;
    \item a \textbf{hypothetical condition} has the imperfect (or present) after \bl{ða}, and the conditional \bl{vel} + infinitive in the main clause;
    \item a \textbf{past counterfactual} has the simple imperfect subjunctive after \bl{ða}, and \bl{volois} + infinitive, or a modal in the imperfect subjunctive, in the main clause.
\end{itemize}
    \glsen{Ða cay ploy je posc, nos voum remanir cas nos.}{ða ke plɔj ʒe pɔx nɔz vum ˌre.maˈnɪr kaz nɔz}{if fall rain {day.of} after {\First\Pl} {will.\First\Pl} {remain-\Inf} {at.home.of} {\First\Pl.\Dsj}}{If it rains tomorrow, we'll stay at home.}
    \glsen{Ða jo saveus, jo volois ig venir.}{ða ʒo saˈvauz ʒo voˈlɔjz aj veˈnɪr}{if {\First\Sg} {know-\Ipfv.\Sbjv} {\First\Sg} {will-\Ipfv.\Sbjv} there {come-\Inf}}{If I had known, I would have come.}
`Unless' is \bl{preðauc}, with the subjunctive.
    \glsen{Preðauc tu quauc, tu kiglara y convey.}{preˈðok ti kwok ti ˌkaj.laˈra i kɔnˈvi}{unless {\Second\Sg} {hurry.\Sbjv} {\Second\Sg} {miss-\Fut} {\Def} train}{Unless you hurry, you'll miss the train.}
`If only' is \bl{ða pog} with the subjunctive (Section~\ref{sec:sbjv-main}). A conditional clause usually comes first, but it may also follow the main clause (\bl{Skig aular ag quarf ða cay ploy borjonnant} `Don't walk to the station if it looks like rain').

\section{Participial and absolute clauses}\label{sec:participial-clauses}
A present participle phrase introduced by \bl{ne} `in' expresses an action going on at the time of the main clause, `while doing'. Its subject is understood to be that of the main clause.
    \glsen{Ne repaðr faint, j'encontrau un veil amig.}{ne reˈpa.ðr̩ fent ˌʒɛn.kɔnˈtro ɪn vil aˈmaj}{in homecoming {do-\Prs.\Ptcp} {\First\Sg=meet-\Pst} {\Indf} old friend}{Walking home, I met an old friend.}
A completed action is expressed with the perfect participle, \bl{ayent} with the object and past participle for transitive verbs and \bl{stant} with the past participle for intransitive and passive ones: \bl{stant marið con un nobr\ldots} `having married a noblewoman\ldots', and after a preposition \bl{lon stant denuð sull'occisour meðes} `without killing himself' (literally `without having become his own killer').

A participle phrase with its own subject forms an absolute construction. It may stand alone, as in \bl{Alcun affer ancour remanent, y raðunanç se closisceu all'hour nouf e demeg} `There being no further business, the meeting closed at half past nine', or be introduced by \bl{ant} `with', as in \bl{l'aziau calm gelað, ant y pousson dy cour se quinant n'eç oregl} `he stood frozen, his pulse thudding in his ears'.

\section{Reported speech}\label{sec:reported-speech}
Direct speech is followed by the verb of saying and its subject, inverted, since the quotation counts as the first constituent of the clause (Section~\ref{sec:v2}): \bl{``\ldots{} donc nel attendr,'' avars la} ```\ldots{} so please don't try,'' she remarked'; \bl{``Cal som nos faint?'' rou Jovanna} ```What are we doing?'' Jovanna asks'.

Reported speech is a complement clause, usually without \bl{ig} (Section~\ref{sec:complement-clauses}). Pronouns shift to the reporter's point of view, as in English. After a past verb of saying, the tense of the reported clause shifts as follows (Section~\ref{sec:tense-subordinate}):
\begin{itemize}
    \item a situation at the same time as the report is in the imperfect;
    \item an earlier event is in the pluperfect;
    \item a later one is in the future in the past, the imperfect of \bl{volir} with an infinitive (\bl{La scavoy la vole ac yembr je Mart}, above).
\end{itemize}
    \glsen{Le dis le fo ja paut.}{le dɪz le fo ʒa pot}{{\Third\Sg.\M} {say.\Pst} {\Third\Sg.\M} {be.\Pst} already {eat.\Pst.\Ptcp}}{He said he had already eaten.}
When the reporter and the reported speaker are the same, an infinitive is often used instead of a finite clause (\bl{Ðe disrn empley broucar n'un foncçondry} `They said they had jobs at an agency'). Reported questions are described in Section~\ref{sec:indirect-questions}, and reported commands and requests take \bl{peðir a} with an infinitive (Section~\ref{sec:complement-clauses}).
